\clearpage
\thispagestyle{empty}

{\centering
\textbf{
    UNIVERSIT\'E LIBRE DE BRUXELLES\\
    Faculté des Sciences\\
    Département d'Informatique
}

\vfill{}

{\Huge Churn Prediction and Causal Analysis \\ on Telecom Customer Data\\}
\vspace{0.5cm}
{\LARGE Théo Verhelst}

\vfill{}

\includegraphics[width=5cm]{figures/seal.pdf}

\vfill{}

\begin{flushleft}
    {\large
    \textbf{Promotor :}     \hfill{} Mémoire présenté en vue de\\
    Prof. Gianluca Bontempi \hfill{} l’obtention du grade de \\
                            \hfill{} Master en Sciences Informatiques
    }
\end{flushleft}

\enlargethispage{1.5cm}
\vfill{}

\textbf{Academic year 2018~-~2019}\par
}

\cleartorecto

\thispagestyle{plain}
\begin{vplace}
\begin{abstract}

Telecommunication companies are evolving in a highly competitive market where
attracting new customers is more expensive than retaining existing ones.
Retention campaigns can be used to prevent customer churn, but this
requires efficient churn prediction models. Such prediction is a difficult
problem, involving large amount of data, non-linearity, imbalance and overlap
between churners and non-churners. In this master thesis, we approach this
problem with Orange Belgium customer data. A descriptive analysis of the dataset
is conducted, and predictive modeling of churn is achieved with a random forest
classifier. The large class imbalance between the two classes is handled with
the Easy Ensemble algorithm. We assess the impact of different data
preprocessing techniques such as feature selection and creation of new features.
The directionality of the impact of variables on churn is estimated through a
sensitivity analysis. Also, we explore the application of data-driven causal
inference, which allows to infer causal relationships between variables purely
from observational data. More specifically, a causal Bayesian network, two
methods of Markov blanket inference, two causal filters and a supervised method
are applied. We draw general conclusions on the possible causes of churn,
supported by the prior knowledge of experts at Orange Belgium.

\end{abstract}
\vfill{}
\begin{otherlanguage}{french}
\begin{abstract}

Les compagnies de télécommunication évoluent dans un marché hautement compétitif
où attirer de nouveaux clients est plus coûteux que retenir les clients déjà
présents. Des campagnes de rétention peuvent être utilisées pour réduire la
résiliation des clients, mais cela nécessite des modèles de prédiction
efficaces. La prédiction de résiliation est un problème difficile, impliquant de
grands quantités de données, des relations non-linéaires, des classes
non-équilibrées et avec une large superposition inter-classe. Dans ce mémoire de
master, nous abordons le problème de prédiction de résiliation avec des données
client de Orange Belgium. Une analyse descriptive du jeu de données est
effectuée, et une modélisation prédictive de la résiliation est obtenue en
utilisant un classificateur random forest. Le non-équilibrage entre les classes
est pris en compte avec l'algorithme Easy Ensemble. Nous évaluons l'impact de
différentes techniques de prétraitement de données, telles que la sélection de
variables et la création de nouvelles variables. La directionalité de l'impact
des variables sur la résiliation est déduite avec une analyse de sensibilité.
Nous explorons également l'utilisation de l'inférence causale, qui permet de
comprendre les liens de causalité entre différentes variables à partir de
données d'observation. Plus spécifiquement, nous utilisons un réseau bayésien
causal, une méthode d'inférence de couverture de Markov, deux filtres causaux et
une méthode supervisée. Nous en tirons des conclusions générales sur les
possibles causes de la résiliation, en prenant en compte les connaissances à
priori d'experts chez Orange Belgium

\end{abstract}
\end{otherlanguage}
\end{vplace}


\clearpage
\chapter*{Acknowledgments}

First and foremost, I would like to thank Prof. Gianluca Bontempi for allowing
me to work with him on this subject and for his unfailing support and
dedication. It is only after these two years that I am able to realize how much
I learned on scientific research and communication with him.

In addition, I am grateful to Olivier Caelen, Jean-Christophe Dewitte, and the
rest of the data science team at Orange for introducing me to their work, for
providing all the resources necessary for this thesis, and for the friendly
atmosphere in the Orange office. In particular, the work of Pierre Brogniet on
churn prediction during his internship brought very valuable and complementary
knowledge to my work.

My special thanks go to my parents, my partner and the rest of my family for
their love and support during these two years and before.

I would like to thank my fellow students at the ULB for the mutual support and
the uncountable science-enthusiastic discussions, fostering my motivation to
pursue scientific research.

Lastly, I would like to express my gratitude to Prof. Tom Lenaerts and Prof.
Bernard Manderick for accepting to evaluate my thesis. Moreover, Prof. Tom
Lenaerts also allowed me to take part in an Erasmus program at Southampton
University, which played a significant role in my academic formation.

\clearpage

\tableofcontents*
%\listoffigures
%\listoftables
